\documentclass{article}
\usepackage{/globals/de}
\usepackage{/globals/math}
\begin{document}
\section{Ebene Geometrie} 
Die Ebene Geometrie ist die vereinfachte Form der Analytischen Geometrie, in einer zweidimensionalen Ebene.

\subsection{Punkte}
Ein Punkt ist ein Paar an zwei reellen Zahlen, eine $x$-Koordinate und eine $y$-Koordinate.
\[
 P = (x;y)
\]
Der Abstand zwischen zwei Punkten $P_1 = (x_1;y_1)$ und $P_2 = (x_2;y_2)$ basiert auf den Satz vom Pythagoras
\[
 d = \sqrt{{\Delta x}^2 + {\Delta y}^2} = \sqrt{(x_2-x_1)^2 + (y_2-y_1)^2}
\]

\subsection{Geraden}
Eine Gerade ist eine unendlich lange Linie in der Ebene. Ein endliches Teil einer Gerade ist ein Segment. Im Zweidimensionalen kann diese durch eine einfache lineare Funktion beschrieben werden.
\[
 y = mx+n \dtext{mit} m = \frac{\Delta x}{\Delta y}
\]
Hier ist die Steigung $m=\Delta x / \Delta y=\tan{\alpha}$. $n$ kann durch eine Punktprobe mit $m$ gefunden werden oder als $n = y_1-mx_1 = y_2-mx_2$.

Gibt es zwei Geraden mit den Steigungen $m_1$ und $m_2$, so sind diese parallel wenn $m_1=m_2$ und senkrecht wenn $m_1*m_2=-1$.

\subsection{Dreiecke}
Drei Punkte einer Ebene können verbunden werden um ein Dreieck zu bilden, solange sie nicht alle auf einer Gerade liegen.

Die Drei Punkte werden \idr $A$, $B$, $C$ genannt. 
An jedem der Punkte liegt ein Innenwinkel, \idr benannt nach dem dazugehörigen griechischen Buchstaben $\alpha$, $\beta$, $\gamma$. Zwischen den Punkten liegen Segmente, benannt nach den kleinbuchstaben der gegenüberliegenden Punkte $a$, $b$, $c$. Für alle Dreiecke gibt
\begin{mathdescription}
 \mathitem{Innenwinkel} \alpha + \beta + \gamma = 180\dgr \\
 \mathitem{Umfang} U = a+b+c \\
 \mathitem{Flächeninhalt} A = \frac{1}{2} b * h
\end{mathdescription}
Es gibt drei relevante Arten an Dreiecken
\begin{description}
 \item[Gleichseitig] Alle Seiten sind gleich lang
 \item[Gleichschenklig] Zwei Seiten sind gleich lang
 \item[Rechtwinklig] Es gibt einen $90\dgr$ Winkel. In einem Rechtwinklig gilt immer
 \begin{mathdescription}
 \mathitem{Pythagoras} a^2+b^2=c^2 \\
 \mathitem{Sinus} \sin() = \frac{\text{Gegenkathete}}{\text{Hypothenuse}} \\
 \mathitem{Kosinus} \cos() = \frac{\text{Ankathete}}{\text{Hypothenuse}} \\
 \mathitem{Tangens} \tan() = \frac{\text{Gegenkathete}}{\text{Ankathete}}
\end{mathdescription}
\end{description}

\subsection{Vierecke}
Vier Punkte, $A$ is $D$, machen ein Viereck. Wenn alle Winkel $90\dgr$ groß sind, ist es ein Rechteck. Wenn alle Seiten gleich lang sind, ist es ein Quadrat.
\[
 U = 2a + 2b \dtext{und} A = a*b
\]

\subsection{Kreise}
Ein Mittelpunkt und ein Radius $r$ beschreiben einen Kreis.
\[
 d=2r \ddtext{und} U = 2 \pi r \dtext{und} \pi r^2
\]
Wird sich ein Punkt $P$ im Winkel $\gamma$ von einem Einheitskreis (Mittelpunkt $M=(0;0), r=1$) angeguckt, so ist der Punkt bei
\[
 P = (\cos{\gamma};\sin{\gamma})
\]

\subsection{Winkel}
Winkel können im Gradmaß und im Bogenmaß gemessen werden; eine Umdrehung sind $360\dgr$ in Gradmaß und $2\pi$ in Bogenmaß. Somit wird mit einer Multiplikation von oder Division von $180\dgr / \pi$ umgewandelt.
\[
 \alpha_B = \alpha_G * \frac{\pi}{180\dgr}
 \dtext{und}
 \alpha_G = \alpha_B * \frac{180\dgr}{\pi}
\]
\end{document}